\scanpage{9}
\mspage{4}
qui conservent resp. renversent l'orientation, i.e.
\[
\operatorname{Hom}_{\mathrm{isot}}^+(X,Y)\quad\text{et}\quad\operatorname{Hom}_{\mathrm{isot}}^-(X,Y).
\]
\deleted{Supposons maintenant}

\uline{Scholie} Par la composition évidente des $\operatorname{Hom}_{\mathrm{isot}}^+(X,Y)$, les surfaces topologiques orientées anabéliennes (i.e. de t.f. et anabéliennes) \deleted{anabéliennes} forment une catégorie, la ``catégorie isotopique orientée'' des surfaces orientées anabéliennes. \deleted{La \unsure{catégorie} $X$ et $Y$ deux surfaces \unsure{anabéliennes}, \unsure{le de}} \inserted{Il y a équivalence entre la catégorie-groupoïde} des triplets $(X,Y,f)$ de deux surfaces \unsure{top.} anabéliennes $X,Y$ et d'un $f\in\operatorname{Hom}_{\mathrm{isot}}^+(X,Y)$, et celle des quadruplets $(X,Y,Y',f_0)$, où $X$ et $Y$ sont itou, $Y'$ un rev. étale fini de $Y$, et $f_0$ est une classe d'isotopie d'un isom. orienté de $X$ et $Y'$, i.e. $f_0\in\operatorname{Isom}_{\mathrm{isot}}^+(X,Y')$.

Posons maintenant
\margin{Pour la suite, on laisse tomber le $+$ dans les notations -- tous les morphismes étales envisagés conservent l'orientation.}
\[
(12)\quad
\widehat{\operatorname{Hom}}^{+}_{\mathrm{isotop}}(X,Y)
=\overbrace{\operatorname{Hom}^{+}_{\mathrm{isot}}(X,Y)\mathop{\wedge}_{\mathcal{T}_X}
\underbrace{\widehat{\mathcal{T}}_X}_{\text{complété profini}}}^{\substack{\text{muni de la topologie}\\\text{somme d'espaces compacts}}}
\]
\margin{[\unsure{Somme} \inserted{\unsure{finie}} d'ailleurs -- donc $\widehat{\operatorname{Hom}}$ est cpt\ldots]}
On définit une composition
\[
(13)\quad\left\{
\begin{aligned}
\widehat{\operatorname{Hom}}_{\mathrm{isot}}(X,\unsure{Y|Z})\times
\widehat{\operatorname{Hom}}_{\mathrm{isot}}(\unsure{X|Y},Z)
&\mathbin{\unsure{-|\longrightarrow}}\widehat{\operatorname{Hom}}_{\mathrm{isot}}(X,Z)\\
\text{contin}\qquad (u,v)&\longmapsto v\circ u
\end{aligned}\right.
\]
\margin{8}
